\section{分类训练与限时训练}

本章按知识模块分类编排练习题，并在章末提供限时模拟卷。

\subsection{力学}

\subsubsection{基础巩固}

\textbf{1.} 一物体做匀加速直线运动，初速度为 $2\,\text{m/s}$，加速度为 $0.5\,\text{m/s}^2$。求：
\begin{enumerate}
  \item 第 $4\,\text{s}$ 末的速度
  \item 前 $4\,\text{s}$ 内的位移
\end{enumerate}

\textbf{2.} 质量为 $5\,\text{kg}$ 的物体静止于水平面上，受到水平拉力 $F = 20\,\text{N}$ 的作用，物体与水平面间的动摩擦因数 $\mu = 0.3$，求物体的加速度。

\textbf{3.} 一物体从 $45\,\text{m}$ 高处自由下落，不计空气阻力，$g = 10\,\text{m/s}^2$。求：
\begin{enumerate}
  \item 落地时间
  \item 落地速度
\end{enumerate}

\subsubsection{能力提升}

\textbf{4.} 如图所示，$A$、$B$ 两物体用轻绳相连跨过光滑定滑轮，$m_A = 2\,\text{kg}$，$m_B = 1\,\text{kg}$，$A$ 放在水平桌面上，$A$ 与桌面间的动摩擦因数 $\mu = 0.2$。释放系统后，求：
\begin{enumerate}
  \item 系统的加速度大小
  \item 绳中的拉力大小
\end{enumerate}

\textbf{5.} 一质量为 $m$ 的小球用长为 $L$ 的轻绳悬挂于 $O$ 点。将小球拉至与 $O$ 点等高的 $A$ 点后由静止释放，不计空气阻力。求：
\begin{enumerate}
  \item 小球运动到最低点 $B$ 时的速度大小
  \item 小球运动到最低点 $B$ 时绳中的拉力大小
  \item 小球摆到右侧最高点 $C$ 时，$OC$ 与竖直方向的夹角
\end{enumerate}

\subsubsection{挑战冲刺}

\textbf{6.} 光滑水平面上，一质量为 $m$ 的滑块以速度 $v_0$ 滑上一质量为 $M$ 的静止长木板，滑块与木板间的动摩擦因数为 $\mu$，木板足够长。求：
\begin{enumerate}
  \item 滑块和木板的共同速度
  \item 滑块相对木板滑行的距离
  \item 系统产生的热量
\end{enumerate}

\textbf{7.} 在光滑水平面上，$A$、$B$ 两球发生弹性正碰。$m_A = 2m$，$m_B = m$，$A$ 球以初速度 $v_0$ 碰撞静止的 $B$ 球。求碰后 $A$ 球和 $B$ 球的速度。

\subsection{曲线运动与万有引力}

\subsubsection{基础巩固}

\textbf{1.} 一小球以 $v_0 = 10\,\text{m/s}$ 的初速度水平抛出，经 $2\,\text{s}$ 后落地，$g = 10\,\text{m/s}^2$。求：
\begin{enumerate}
  \item 小球下落的高度
  \item 水平射程
  \item 落地时的速度大小
\end{enumerate}

\textbf{2.} 已知地球半径为 $R$，地球表面重力加速度为 $g$。求地球的第一宇宙速度表达式。

\textbf{3.} 一人造卫星绕地球做匀速圆周运动的轨道半径为 $r$，地球质量为 $M$，万有引力常量为 $G$。求卫星的线速度 $v$ 和周期 $T$。

\subsubsection{能力提升}

\textbf{4.} 一质量为 $m$ 的小球在长为 $L$ 的细绳牵引下在竖直面内做圆周运动。求小球在最高点时细绳的最小拉力。

\textbf{5.} 已知月球的质量是地球质量的 $\dfrac{1}{81}$，月球的半径是地球半径的 $\dfrac{1}{4}$。求：
\begin{enumerate}
  \item 月球表面的重力加速度与地球表面重力加速度之比
  \item 在月球上发射卫星的最小速度与地球第一宇宙速度之比
\end{enumerate}

\subsubsection{挑战冲刺}

\textbf{6.} 如图所示，在倾角为 $\theta$ 的斜面顶端以水平速度 $v_0$ 抛出一个小球，小球落在斜面上。求：
\begin{enumerate}
  \item 小球在空中飞行的时间
  \item 小球落到斜面上时的速度方向与水平方向的夹角
\end{enumerate}

\subsection{电场与电路}

\subsubsection{基础巩固}

\textbf{1.} 真空中有两个点电荷，相距 $0.1\,\text{m}$，电荷量分别为 $q_1 = 2\times10^{-8}\,\text{C}$ 和 $q_2 = -3\times10^{-8}\,\text{C}$。求两电荷间的作用力大小和方向。

\textbf{2.} 在匀强电场中，将一电荷量为 $q = 1\times10^{-8}\,\text{C}$ 的正电荷从 $A$ 点移到 $B$ 点，电场力做功 $W = 5\times10^{-6}\,\text{J}$，$A$、$B$ 间距 $0.2\,\text{m}$，沿电场线方向。求：
\begin{enumerate}
  \item $A$、$B$ 间的电势差
  \item 电场强度的大小
\end{enumerate}

\textbf{3.} 如图所示电路，$R_1 = 2\,\Omega$，$R_2 = 3\,\Omega$，$R_3 = 6\,\Omega$，电源电动势 $E = 6\,\text{V}$，内阻 $r = 1\,\Omega$。求：
\begin{enumerate}
  \item 外电路总电阻
  \item 干路电流
  \item 路端电压
\end{enumerate}

\subsubsection{能力提升}

\textbf{4.} 一个平行板电容器，极板面积为 $S$，板间距离为 $d$，充电后两极板间电压为 $U$，断开电源。现保持极板面积不变，将板间距离增大到 $2d$。求：
\begin{enumerate}
  \item 电容变为原来的多少倍
  \item 两极板间的电场强度如何变化
\end{enumerate}

\textbf{5.} 如图所示，电源电动势 $E = 10\,\text{V}$，内阻 $r = 0.5\,\Omega$，$R_1 = 4\,\Omega$，$R_2$ 为滑动变阻器（$0\sim20\,\Omega$）。当 $R_2$ 取何值时，电源输出功率最大？最大输出功率是多少？

\subsubsection{挑战冲刺}

\textbf{6.} 一质量为 $m$、电荷量为 $+q$ 的粒子以初速度 $v_0$ 沿垂直于电场方向射入匀强电场，极板长为 $L$，板间距离为 $d$，极板间电压为 $U$。不计重力，求：
\begin{enumerate}
  \item 粒子在电场中的运动时间
  \item 粒子的偏转量
  \item 粒子射出电场时的速度偏转角
\end{enumerate}

\subsection{磁场与电磁感应}

\subsubsection{基础巩固}

\textbf{1.} 带电粒子以速度 $v$ 垂直射入磁感应强度为 $B$ 的匀强磁场中，粒子电荷量为 $q$、质量为 $m$。求粒子做圆周运动的半径 $r$ 和周期 $T$。

\textbf{2.} 一根长 $L = 0.4\,\text{m}$ 的直导线，通有 $I = 2\,\text{A}$ 的电流，放在磁感应强度 $B = 0.5\,\text{T}$ 的匀强磁场中。求：
\begin{enumerate}
  \item 导线与磁场垂直时受到的安培力
  \item 导线与磁场方向成 $30^\circ$ 角时受到的安培力
\end{enumerate}

\textbf{3.} 一矩形线圈在匀强磁场中绕垂直于磁场的轴匀速转动，产生正弦式交变电流的电动势 $e = 100\sin100\pi t\;\text{(V)}$。求：
\begin{enumerate}
  \item 电动势的最大值和有效值
  \item 交流电的频率
  \item 当 $t = 0.005\,\text{s}$ 时的电动势瞬时值
\end{enumerate}

\subsubsection{能力提升}

\textbf{4.} 如图所示，$MN$ 和 $PQ$ 是两根平行金属导轨，$L = 0.5\,\text{m}$，电阻 $R = 3\,\Omega$，导轨平面与水平面成 $30^\circ$ 角。匀强磁场竖直向上，$B = 1\,\text{T}$。导体棒 $ab$ 质量 $m = 0.1\,\text{kg}$，电阻 $r = 1\,\Omega$，从静止释放，导轨电阻不计。求：
\begin{enumerate}
  \item 导体棒 $ab$ 能达到的最大速度
  \item 当 $v = 2\,\text{m/s}$ 时，导体棒的加速度
\end{enumerate}

\textbf{5.} 一理想变压器原线圈匝数 $n_1 = 1000$，副线圈匝数 $n_2 = 200$，原线圈接 $220\,\text{V}$ 交流电。求：
\begin{enumerate}
  \item 副线圈的输出电压
  \item 若副线圈接 $44\,\Omega$ 的负载，原线圈的输入电流
\end{enumerate}

\subsubsection{挑战冲刺}

\textbf{6.} 在 $xOy$ 平面内，第 I 象限存在匀强磁场，磁感应强度为 $B$，方向垂直纸面向里。一质量为 $m$、电荷量为 $+q$ 的粒子以速度 $v$ 从 $y$ 轴上的 $P(0,a)$ 点沿 $+x$ 方向射入磁场。求：
\begin{enumerate}
  \item 粒子在磁场中的运动半径
  \item 粒子在磁场中运动的时间
  \item 粒子射出磁场时与 $x$ 轴的交点坐标
\end{enumerate}

\subsection{热学（选修 3--3）}

\textbf{1.} 一定质量的理想气体，从状态 $A$（$p_A = 1.0\times10^5\,\text{Pa}$，$V_A = 2.0\,\text{L}$，$T_A = 300\,\text{K}$）经等温压缩到状态 $B$（$V_B = 1.0\,\text{L}$），再经等容升温到状态 $C$（$T_C = 600\,\text{K}$）。求：
\begin{enumerate}
  \item 状态 $B$ 的压强 $p_B$
  \item 状态 $C$ 的压强 $p_C$
\end{enumerate}

\textbf{2.} 如图所示，一气缸内封闭一定质量的理想气体，活塞质量为 $m$，横截面积为 $S$，大气压强为 $p_0$，活塞与气缸壁间无摩擦。求气缸内气体的压强。

\textbf{3.} 一定质量的理想气体从外界吸收 $200\,\text{J}$ 的热量，同时气体对外界做功 $80\,\text{J}$。求：
\begin{enumerate}
  \item 气体内能的变化量
  \item 气体温度升高还是降低
\end{enumerate}

\subsection{光学（选修 3--4）}

\textbf{1.} 一束光从折射率为 $n = 1.5$ 的某种介质射向空气，入射角为 $30^\circ$。求：
\begin{enumerate}
  \item 折射角
  \item 光在介质中的传播速度（$c = 3\times10^8\,\text{m/s}$）
\end{enumerate}

\textbf{2.} 水的折射率为 $n = \frac43$，求光从水射向空气时的全反射临界角。

\textbf{3.} 在杨氏双缝干涉实验中，双缝间距 $d = 0.2\,\text{mm}$，缝到屏的距离 $L = 1\,\text{m}$，所用单色光波长为 $\lambda = 500\,\text{nm}$。求相邻亮纹中心的距离。

\subsection{限时训练}

\subsubsection{限时训练一（基础卷）}

\textbf{时间}：30 分钟 \quad \textbf{满分}：40 分

\textbf{一、单选题（每题 4 分，共 20 分）}

\textbf{1.} 一物体做自由落体运动，取 $g = 10\,\text{m/s}^2$，第 $2\,\text{s}$ 末的速度为
\begin{enumerate}
  \item[A.] $10\,\text{m/s}$ \quad [B.] $15\,\text{m/s}$ \quad [C.] $20\,\text{m/s}$ \quad [D.] $25\,\text{m/s}$
\end{enumerate}

\textbf{2.} 两个点电荷之间的库仑力大小为 $F$，若将它们之间的距离增大为原来的 2 倍，同时使其中一个电荷量减半，则库仑力变为
\begin{enumerate}
  \item[A.] $\dfrac18F$ \quad [B.] $\dfrac14F$ \quad [C.] $\dfrac12F$ \quad [D.] $F$
\end{enumerate}

\textbf{3.} 一正弦式交流电的有效值为 $10\,\text{V}$，其最大值为
\begin{enumerate}
  \item[A.] $5\sqrt{2}\,\text{V}$ \quad [B.] $10\sqrt{2}\,\text{V}$ \quad [C.] $20\,\text{V}$ \quad [D.] $20\sqrt{2}\,\text{V}$
\end{enumerate}

\textbf{4.} 氢原子从 $n=3$ 能级跃迁到 $n=2$ 能级时辐射光子的能量为
\begin{enumerate}
  \item[A.] $1.89\,\text{eV}$ \quad [B.] $3.40\,\text{eV}$ \quad [C.] $10.2\,\text{eV}$ \quad [D.] $12.09\,\text{eV}$
\end{enumerate}

\textbf{5.} 光从空气射入水中，下列说法正确的是
\begin{enumerate}
  \item[A.] 频率变大 \quad [B.] 波长变大 \quad [C.] 速度变小 \quad [D.] 颜色改变
\end{enumerate}

\textbf{二、计算题（20 分）}

\textbf{6.} 一质量为 $2\,\text{kg}$ 的物体在水平面上受到 $10\,\text{N}$ 的水平拉力，物体与水平面间的动摩擦因数 $\mu = 0.2$，$g = 10\,\text{m/s}^2$。求：
\begin{enumerate}
  \item 物体的加速度
  \item 物体从静止开始运动 $4\,\text{s}$ 后的速度
\end{enumerate}

\subsubsection{限时训练二（综合卷）}

\textbf{时间}：45 分钟 \quad \textbf{满分}：50 分

\textbf{一、多选题（每题 5 分，共 15 分）}

\textbf{1.} 关于电场强度和电势，下列说法正确的是
\begin{enumerate}
  \item[A.] 电场强度为零的地方，电势一定为零
  \item[B.] 电势为零的地方，电场强度一定为零
  \item[C.] 电场强度大的地方，电势不一定高
  \item[D.] 沿着电场线方向，电势逐渐降低
\end{enumerate}

\textbf{2.} 关于机械能守恒，下列说法正确的是
\begin{enumerate}
  \item[A.] 做自由落体运动的物体机械能守恒
  \item[B.] 做匀速圆周运动的物体机械能一定守恒
  \item[C.] 物体所受合外力为零时机械能一定守恒
  \item[D.] 只有重力做功时机械能守恒
\end{enumerate}

\textbf{3.} 关于电磁感应，下列说法正确的是
\begin{enumerate}
  \item[A.] 穿过闭合回路的磁通量变化越大，感应电动势越大
  \item[B.] 感应电流的磁场总是阻碍原磁通量的变化
  \item[C.] 导体切割磁感线的速度越大，感应电动势越大
  \item[D.] 感应电流的磁场方向可能与原磁场方向相同
\end{enumerate}

\textbf{二、实验题（15 分）}

\textbf{4.} 在"验证牛顿第二定律"的实验中：
\begin{enumerate}
  \item 平衡摩擦力时，应 \underline{\hspace{3cm}}
  \item 实验中得到一条纸带，相邻计数点间时间为 $0.1\,\text{s}$，各点间距分别为 $x_1=2.00\,\text{cm}$、$x_2=3.50\,\text{cm}$、$x_3=5.00\,\text{cm}$、$x_4=6.50\,\text{cm}$，求加速度 $a$
\end{enumerate}

\textbf{三、计算题（20 分）}

\textbf{5.} 光滑水平面上，$A$ 球以 $6\,\text{m/s}$ 的速度向右运动，与静止的 $B$ 球发生正碰。$m_A = 2\,\text{kg}$，$m_B = 4\,\text{kg}$，碰撞后 $A$ 球以 $2\,\text{m/s}$ 的速度反弹。求：
\begin{enumerate}
  \item $B$ 球碰后的速度
  \item 碰撞过程中系统损失的动能
\end{enumerate}
